\clearpage\begingroup\SourceLeaf{32b}\HistoricalHead{27}{R}{GAUSS AN SCHUMACHER.}
\begin{center}
[5.]\\[.7em]
\textsc{Gauss} an \textsc{Schumacher}. Göttingen, 1. Julius 1836.
\par\medskip\rule{.78\textwidth}{.35pt}\par
{\small [Briefwechsel zwischen \textsc{Gauss} und \textsc{Schumacher} III, Altona 1860, S. 83, 84.]}\par
\rule{.78\textwidth}{.35pt}
\end{center}
\noindent\hspace*{\parindent}Bei \textsc{Olufsens} Wägungen in Ihrem J[ahr] B[uch, 1836, S. 244 ff.] habe ich
den M[ittleren] F[ehler] Einer Wägung \(=0^{\mathrm{P}}{,}119=0^{\mathrm{mgr.}}{,}565\)[\textsuperscript{*})]\begingroup\renewcommand{\thefootnote}{}\footnotetext{[\textsuperscript{*}) Bei \(0^{\mathrm{P}}{,}119\) ist P eine Abkürzung für \textls[120]{partes} und bedeutet die Teile der Skala, auf der Zeiger des Waagebalkens schwingt. Nach \textsc{Olufsen} a. a. O. S. 246 sind 3,37 Teile gleich 16 mgr., was mit der Angabe von \textsc{Gauss} übereinstimmt.]}\endgroup; die gegen-\HistoricalFoot{}{4\textsuperscript{*}}
\clearpage\SourceLeaf{33}\HistoricalHead{28}{L}{PHYSIK. BRIEFWECHSEL.}
\noindent wärtigen 6 würden, obgleich die Belastung nicht einmal halb so gross ist,
etwa \(0^{\mathrm{mgr.}}{,}8\) geben! Ist nun die Pariser Waage so viel besser\textsuperscript{*})\begingroup\renewcommand{\thefootnote}{}\footnotetext{\textsuperscript{*}) Ich setze jedoch nicht in Abrede, dass der M. F. aus so wenigen Wägungen gar nicht mit Sicherheit bestimmt werden kann.}\endgroup, oder handelt
es sich um einen mir noch entgangenen Umstand in der Behandlung? Ich
brauche wohl nicht zu sagen, dass ich die Operation mit äusserster Vorsicht
mache, namentlich
\begin{list}{}{\setlength{\leftmargin}{2.5em}\setlength{\labelwidth}{1.3em}\setlength{\labelsep}{.4em}\setlength{\itemsep}{0pt}\setlength{\parsep}{0pt}\setlength{\topsep}{0pt}}
\item[1)] wird die Waage, oder der Tisch, worauf sie steht, gar nicht von der
Sonne beschienen, oder von irgend einer unregelmässigen Wärmequelle
benachbart.
\item[2)] Abheben und Aufsetzen der Gewichte geschieht sehr vorsichtig mit
einer Gabel, während die Waage gestützt ist.
\item[3)] Hinauf- und Herunterschrauben der Stütze mache ich \textls[120]{sehr} langsam.
\item[4)] Thüren, sowohl der Waage als des Zimmers, natürlich auch Fenster
sind verschlossen.
\item[5)] Ich beobachte in einem \textls[120]{andern} Zimmer mit Fernrohr, wozu ich mir
eine ganz bequeme Einrichtung getroffen habe (Entfernung etwa 20
p[ariser] Fuss).
\end{list}
\par\noindent\hspace*{\parindent}Bis jetzt sehe ich nun keinen andern Erklärungsgrund als folgende:
\begin{list}{}{\setlength{\leftmargin}{2.5em}\setlength{\labelwidth}{1.3em}\setlength{\labelsep}{.4em}\setlength{\itemsep}{0pt}\setlength{\parsep}{0pt}\setlength{\topsep}{0pt}}
\item[1)] Die nach der Aushebung wieder gesenkte Waage kommt nicht in absoluter
Schärfe wieder auf die nämliche Stelle der Agatplatte; dies
kann freilich nur insofern einen Unterschied machen, als die Agatfläche
kein \textls[120]{vollkommenes Planum} ist (selbst eine Inclination,
wenn jene wirklich ein Planum ist, könnte \textls[120]{keinen} Unterschied
machen), und \textls[120]{vollkommen} ist freilich kein Menschenwerk.
\item[2)] Mehr Wahrscheinlichkeit hat aber wohl eine zweite Ursache: die
beiden Tragstücke setzen sich beim Niederlassen nicht in absoluter
Schärfe wieder auf die nemlichen Stellen der beiden äussern Schneiden.
Dies macht aber einen Unterschied, wenn die Schneiden nicht absolut
genau gerade, mit der mittlern Schneide parallele Linien sind, und
selbst eine \textls[120]{sehr kleine} Abweichung in dieser Beziehung könnte doch
wohl im vorliegenden Fall merklich werden. Das Mittel, welches (wie
schon gesagt, nach \textsc{Webers} Autorität) die Künstler zu diesem Zweck
\clearpage\SourceLeaf{34}\HistoricalHead{29}{R}{GAUSS AN SCHUMACHER.}
\noindent anwenden, ist doch wahrlich nur ein sehr rohes, und man muss sich
daher fast wundern, dass demungeachtet die Waagen noch so viel
leisten. Ich habe mir die Sache etwas hin und her überlegt, und
möchte für das durchschlagendste Prüfungsmittel halten, wenn man
die Einrichtung träfe, dass jene Drahtstücke \textls[120]{behuf der Prüfung}
zuerst einen beträchtlichen Spielraum längs der Schneide hätten, indem
man sie z. B. \textls[120]{einmahl} weit nach vorne, und dann auch weit nach
hinten zu aufliegen lassen könnte, und so lange corrigirte, bis bei
beiden Auflagen kein merklich verschiedenes Resultat käme, nachher
aber, für den wirklichen Gebrauch, den Spielraum in sehr enge
Grenzen einschlösse. Was ist Ihre Meinung darüber?\quad .\ .\ .\ .\ .\ .\ .\ .\ .
\end{list}
\HistoricalHead{29}{R}{GAUSS AN SCHUMACHER.}
\clearpage\endgroup
